% Prospective (not yet executed) research endpoint and E0 pilot, 12 Nov 2026.
% Preserve M5/M6 historical 500-sweep results; do not re-label them.
\section{Prospective research decision: bits worth storing}
\label{sec:prospective-design}

The completed M5/M6 evidence establishes that a six-bit
floating-point representation is not uniformly preferable
to a six-bit integer representation across weighted
Max-Cut workloads. Static local-field and switching-probability
fidelity rankings reverse between the two held-out graphs,
while the finite-budget quality contrasts are partly
inconclusive. This motivates a narrower, explicitly
hardware-oriented \emph{next-stage} decision, distinct
from the already completed experimental claims:
\begin{quote}\itshape
For digitally stored Ising couplings feeding STT-MRAM
p-bits, how many coefficient payload bits, if any,
can low-order floating point save relative to integer
encoding at comparable original-weight Max-Cut
quality, and does realistic p-bit variation change
that minimum useful precision?
\end{quote}
For a fixed graph family, response condition, solver
budget, pairing and \emph{prospectively} agreed
quality criterion $\Delta$, denote the smallest
eligible FP or integer payload width as
$b^*_{\mathrm{FP}}$ and $b^*_{\mathrm{INT}}$.
The proposed primary endpoint is
\[
  k=b^*_{\mathrm{INT}}-b^*_{\mathrm{FP}}
  \quad\text{bits saved per stored coupling}.
\]
If a width is not certified at the sampled precision
and statistical resolution, $b^*$ is reported as
unresolved rather than assigned by interpolation.
Previously observed non-monotonic quality outcomes
also preclude assuming that every larger width
must perform better without evidence. The
proposed tolerance $\Delta=0.5$ percentage points
of added original-optimum gap against same-response
FP64 is a candidate for supervisor agreement;
the completed M5/M6 trials used a different
predeclared pair of 5\% absolute and $+1$ pp
extra-gap screening bounds. They must not be
retroactively re-scored as if the new criterion
had been frozen beforehand.

At 100 spins, symmetric triangular dense
storage has $N(N-1)/2=4950$ unique off-diagonal
positions, so a positive $k$ would represent
$4950k$ \emph{nominal} payload bits saved,
before index and scale metadata. An edge-list
implementation for the 200-nonzero-edge M6
instance would instead save $200k$ bits.
Whether these payload differences compensate
for FP decode/accumulate logic depends on
the memory architecture and updated equal-quality
synthesis; they are not measured energy savings.

\subsection{First gate completed: FP64 calibration on CAL graphs}
The bounded E0-A reference solver audit is now
\emph{complete} on the two already-viewed CAL
graphs, under a source-frozen protocol.
Section~\ref{sec:e0-calibration} gives the
paired results and time-resolved acceptance
diagnostics: the geometric schedule did not
provide a statistically established final
quality improvement over legacy linear at the
same 50,000-update budget on either graph.
Neither schedule reached the exploratory
1\% mean optimum-gap aspiration. The existing
linear 50,000-update condition is retained
as a \emph{finite-budget reference}, not as
an assertion of solver convergence.

The next protocol should therefore focus on
\emph{minimum certified coupling width},
rather than on further ad hoc solver annealing
tuning. Two decisions require supervisor
agreement \emph{before} new precision trials:
(i) fix the same 50,000-update linear solver
budget for reference FP64 and all formats,
and (ii) use a proposed upper 95\% paired
confidence-bound tolerance of $+0.5$
\emph{percentage points} of extra original-weight
optimum gap relative to FP64 under the same
graph, response, paired seeds and budget.
For example, an FP64 mean gap of 4.0\% and
a low-bit mean gap of 4.5\% imply a 0.5-pp
\emph{mean} difference, but eligibility
requires an upper confidence bound at most
$+0.5$ pp; the means alone do not suffice.

The completed M5 $d=0.8$ study illustrates
the potential width decision without
retroactively changing its protocol:
FP6 had an extra-gap upper 95\% confidence
bound of $+0.241$ pp and INT6 of
$+1.287$ pp. A new, separately
preregistered study must determine
$b^*_{\mathrm{FP}}$ and
$b^*_{\mathrm{INT}}$ across the
specified candidate widths. It must
not infer that INT7, INT8, or a lower
FP width meets the tolerance
without evidence. If the predetermined
sample cap leaves a boundary-straddling
confidence interval, the minimum
useful width will be reported
as \emph{unresolved}, not selected
post hoc.

Subsequently, any minimum-width study,
independently sampled device gain/offset
populations, and memory-inclusive
equal-quality cost comparison will require
separately frozen protocols. No
floating-point storage advantage, device
masking effect, solver-convergence target
or post-calibration physical-device statistic
is asserted by the prospective plan alone.
